\chapter{Нормированная мера рождения ячеек и граница темпа роста}
\label{ch:v10-cell-birth-normalized-transition-measure-growth-rate-origin}

\section{Проверяемая гипотеза}

Предыдущий гейт вывел спектральный ход относительно безразмерной координаты
$\zeta$, но не определил вероятности переходов $N\to N+1$. Настоящий гейт
проверяет минимальную последовательную меру рождения: достаточно ли
нормировать вакуумный вес $e^{-S_{\rm vac}}$, чтобы одновременно получить
закон роста и ранее условную амплитуду
\begin{equation}
 h_{\rm vac}=\frac{e^{-S_{\rm vac}}}{\sqrt{8\pi}}.
 \label{eq:v10-birth-target-amplitude}
\end{equation}

Проверка не отождествляет порядковый номер рождения с физическим временем.
Именно возможность такого отождествления должна следовать из меры, а не
предполагаться заранее.

\section{Минимальная нормированная мера одного шага}

Обозначим
\begin{equation}
 x:=e^{-S_{\rm vac}}>0.
 \label{eq:v10-birth-weight}
\end{equation}
Пусть относительные веса исходов ``нет рождения'' и ``одна новая ячейка''
равны соответственно $1$ и $x$. Единственная нормировка этих двух весов
даёт
\begin{equation}
 p_0=\frac1{1+x},
 \qquad
 p_1=\frac{x}{1+x},
 \qquad
 p_0+p_1=1,
 \qquad
 \frac{p_1}{p_0}=x.
 \label{eq:v10-birth-normalized-probabilities}
\end{equation}
Минимальная матрица перехода
\begin{equation}
 T_x=
 \begin{pmatrix}
 p_0&p_1\\0&1
 \end{pmatrix}
 \label{eq:v10-birth-transition-matrix}
\end{equation}
является построчно стохастической. Это честная положительная мера, но она
условна на выборе именно двух исходов и отношения весов $x$.

Если каждая существующая ячейка независимо сохраняется и с вероятностью
$p_1$ производит одну дополнительную ячейку, средний множитель числа ячеек
за один порядковый шаг равен
\begin{equation}
 m(x)=1+p_1=\frac{1+2x}{1+x}.
 \label{eq:v10-birth-mean-multiplier}
\end{equation}
Поэтому
\begin{equation}
 \mathbb E N_n=N_0m(x)^n,
 \qquad
 \Delta\zeta_{\rm step}(x)
 =\frac13\log m(x)
 =\frac13\log\frac{1+2x}{1+x}.
 \label{eq:v10-birth-dimensionless-growth}
\end{equation}
Так нормировка действительно создаёт безразмерный стохастический закон
роста.

\section{Несовпадение с условной вакуумной амплитудой}

В слабовесовом пределе точные линейные коэффициенты равны
\begin{equation}
 \left.\frac{d\Delta\zeta_{\rm step}}{dx}\right|_{x=0}
 =\frac13,
 \qquad
 \left.\frac{dh_{\rm vac}}{dx}\right|_{x=0}
 =\frac1{\sqrt{8\pi}}.
 \label{eq:v10-birth-slope-comparison}
\end{equation}
Причём
\begin{equation}
 \frac13-\frac1{\sqrt{8\pi}}>0.
 \label{eq:v10-birth-prefactor-mismatch}
\end{equation}
Следовательно, функции
$\Delta\zeta_{\rm step}(x)$ и $x/\sqrt{8\pi}$ не тождественны. Это не
опровергает формулу для $h_{\rm vac}$: оно доказывает более узкое
утверждение, что двухисходная нормировка весов сама по себе её не выводит.
Для совпадения требуется дополнительная длительность шага, другой набор
исходов или новая мера.

\section{Непрерывные ожидания и свободная скорость часов}

Даже если перейти к непрерывному чистому процессу рождения с полной
скоростью $r_N=\gamma N$, плотность времени ожидания
\begin{equation}
 f_N(t)=r_Ne^{-r_Nt}
 \label{eq:v10-birth-waiting-density}
\end{equation}
нормирована при любом $r_N>0$:
\begin{equation}
 \int_0^\infty f_N(t)\,dt=1,
 \qquad
 \mathbb E(t\mid N)=\frac1{r_N}.
 \label{eq:v10-birth-waiting-normalization}
\end{equation}
Преобразование
\begin{equation}
 \gamma\longmapsto c\gamma,
 \qquad
 t\longmapsto\frac{t}{c}
 \label{eq:v10-birth-clock-rescaling}
\end{equation}
сохраняет $e^{-\gamma Nt}$. В логарифмических координатах карта
$(\log\gamma,\log t)\mapsto\log\gamma+log t$ имеет одномерное ядро.
Нормировка вероятности поэтому не фиксирует физическую единицу времени.

\begin{theorem}[Минимальная мера рождения нормируется]
Отношение весов $(1,x)$ однозначно задаёт вероятности
$p_0=1/(1+x)$ и $p_1=x/(1+x)$. Независимый процесс рождения имеет точный
средний множитель $(1+2x)/(1+x)$ и безразмерный прирост
$\Delta\zeta_{\rm step}=\frac13\log((1+2x)/(1+x))$.
\end{theorem}

\begin{nogocriterion}[Нормировка не создаёт часы]
Нормированная дискретная мера не воспроизводит коэффициент
$1/\sqrt{8\pi}$ без дополнительной калибровки шага. Нормированная
непрерывная мера существует при любом положительном $\gamma$ и сохраняет
орбиту $(\gamma,t)\mapsto(c\gamma,t/c)$. Поэтому ни вероятность, ни
физический темп роста не следуют из одной нормировки $e^{-S_{\rm vac}}$.
\end{nogocriterion}

\begin{protocol}[Статус меры рождения]\begin{enumerate}
\item положительная двухисходная мера: \textbf{построена};
\item нормировка вероятностей: \textbf{$1$};
\item точный безразмерный закон среднего роста: \textbf{получен};
\item архитектура минимальной меры: \textbf{$8/8$};
\item общий реестр происхождения: \textbf{$2/4$};
\item коэффициент $1/\sqrt{8\pi}$ из нормировки: \textbf{$0/1$};
\item физическая скорость часов $\gamma$: \textbf{$0/1$};
\item физическая космологическая амплитуда: \textbf{не выведена};
\item следующий гейт: \textbf{общий родитель энергии часов и скорости
рождения ячеек}.
\end{enumerate}\end{protocol}

Машинный сертификат сохранён в
\path{s2t_v10_cell_birth_normalized_transition_measure_growth_rate_origin_gate_results.json}.

\begin{thebibliography}{9}
\bibitem{v10-birth-rideout-sorkin}
D.~P.~Rideout, R.~D.~Sorkin,
\emph{A Classical Sequential Growth Dynamics for Causal Sets},
Physical Review D 61 (2000), 024002.

\bibitem{v10-birth-varadarajan-rideout}
M.~Varadarajan, D.~Rideout,
\emph{A General Solution for Classical Sequential Growth Dynamics of
Causal Sets},
Physical Review D 73 (2006), 104021.
\end{thebibliography}